\documentclass{beamer}
\usepackage[utf8]{inputenc}

\usetheme{Madrid}
\usecolortheme{default}
\usepackage{amsmath,amssymb,amsfonts,amsthm}
\usepackage{txfonts}
\usepackage{tkz-euclide}
\usepackage{listings}
\usepackage{adjustbox}
\usepackage{array}
\usepackage{tabularx}
\usepackage{gvv}
\usepackage{lmodern}
\usepackage{circuitikz}
\usepackage{lmodern}
\usepackage{multicol}
\usepackage{tikz}
\usepackage{graphicx}

\setbeamertemplate{page number in head/foot}[totalframumber]

\usepackage{tcolorbox}
\tcbuselibrary{minted,breakable,xparse,skins}



\definecolor{bg}{gray}{0.95}
\DeclareTCBListing{mintedbox}{O{}m!O{}}{%
  breakable=true,
  listing engine=minted,
  listing only,
  minted language=#2,
  minted style=default,
  minted options={%
    linenos,
    gobble=0,
    breaklines=true,
    breakafter=,,
    fontsize=\small,
    numbersep=8pt,
    #1},
  boxsep=0pt,
  left skip=0pt,
  right skip=0pt,
  left=25pt,
  right=0pt,
  top=3pt,
  bottom=3pt,
  arc=5pt,
  leftrule=0pt,
  rightrule=0pt,
  bottomrule=2pt,
  toprule=2pt,
  colback=bg,
  colframe=orange!70,
  enhanced,
  overlay={%
    \begin{tcbclipinterior}
    \fill[orange!20!white] (frame.south west) rectangle ([xshift=20pt]frame.north west);
    \end{tcbclipinterior}},
  #3,
}
\lstset{
    language=C,
    basicstyle=\ttfamily\small,
    keywordstyle=\color{blue},
    stringstyle=\color{orange},
    commentstyle=\color{green!60!black},
    numbers=left,
    numberstyle=\tiny\color{gray},
    breaklines=true,
    showstringspaces=false,
}
%------------------------------------------------------------
%This block of code defines the information to appear in the
%Title page
\title %optional
{12.640}
\date{October 11,2025}
%\subtitle{A short story}

\author % (optional)
{EE25BTECH11002 - Achat Parth Kalpesh}



\begin{document}


\frame{\titlepage}

\begin{frame}{Question}
$\vec{P}$ and $\vec{R}$ are Jacobian matrices for two different geophysical inverse problems. If their generalized inverses are written as $\vec{P}^{-g} = \brak{\vec{P}^T \vec{P}}^{-1} \vec{P}^T$ and $\vec{R}^{-g} = \vec{R}^T\brak{\vec{R} \vec{R}^T}^{-1}$, then
\begin{enumerate}
    \item $\vec{P}$ represents an overdetermined problem and $\vec{R}$ represents an underdetermined problem
    \item $\vec{P}$ represents an underdetermined problem and $\vec{R}$ represents an overdetermined problem
    \item Both $\vec{P}$ and $\vec{R}$ represent overdetermined problems
    \item Both $\vec{P}$ and $\vec{R}$ represent underdetermined problems
\end{enumerate}
\end{frame}

\begin{frame}{Solution}
Given are two Jacobian matrices $\vec{P}$ and $\vec{R}$ with generalized inverses:
\begin{align}
    \vec{P}^{-g} &= \brak{\vec{P}^T \vec{P}}^{-1} \vec{P}^T \\
    \vec{R}^{-g} &= \vec{R}^T \brak{\vec{R} \vec{R}^T}^{-1}
\end{align}
Let $\vec{P} \in \mathbb{R}^{m \times n}$ and $\vec{R} \in \mathbb{R}^{m \times n}$
\begin{align}
    \vec{P}^{-g} = \brak{\vec{P}^T \vec{P}}^{-1} \vec{P}^T\\
    \vec{P}^T \vec{P} \in \mathbb{R}^{n \times n}
\end{align}
\end{frame}



\begin{frame}{Solution}
For the inverse to exist
\begin{align}
   \det\brak{\vec{P}^T \vec{P}} \neq 0 \\
   \implies \text{rank}\brak{\vec{P}} = n \leq m
\end{align}
This implies $m \geq n$, i.e., more equations than unknowns. Hence, $\vec{P}$ is tall and the system is overdetermined
\begin{align}
    \vec{R}^{-g} = \vec{R}^T \brak{\vec{R} \vec{R}^T}^{-1}\\
    \vec{R} \vec{R}^T \in \mathbb{R}^{m \times m}
\end{align}
For the inverse to exist
\begin{align}
    \det\brak{\vec{R} \vec{R}^T} \neq 0 \\
    \implies \text{rank}\brak{\vec{R}} = m \leq n
\end{align}
\end{frame}

\begin{frame}{Solution}
This implies $n \geq m$, i.e., more unknowns than equations. Hence, $\vec{R}$ is wide and the system is underdetermined\\
$\vec{P}$ represents an overdetermined system and $\vec{R}$ represents an underdetermined system
\end{frame}


\end{document}